\section{Distribuciones y Compatibilidad}

\subsection{El esquema de versiones}

ROS~2 sigue un ciclo de lanzamiento sincronizado con Ubuntu LTS. El motivo es pragmático: Ubuntu LTS (Long Term Support) garantiza cinco años de soporte de seguridad para sus paquetes del sistema, y los paquetes de ROS~2 dependen fuertemente de las versiones del sistema: Python, CMake, OpenSSL, las bibliotecas de DDS, y los compiladores. Construir ROS~2 sobre una base de sistema estable reduce drásticamente los problemas de dependencias entre usuarios~\cite{ros2releases}.

Las distribuciones LTS de ROS~2 se lanzan junto con las versiones LTS de Ubuntu (cada dos años, en abril) y reciben mantenimiento por cinco años. Las distribuciones no-LTS (a veces llamadas STD, Standard) se lanzan en los años intermedios, tienen soporte de un año, y sirven como terreno de prueba para características nuevas que eventualmente llegan a la siguiente LTS. Para proyectos de investigación de largo plazo, las distribuciones LTS son siempre la elección correcta.

El nombre de cada distribución sigue una convención de dos palabras aliteradas con nombres de animales en orden alfabético: Ardent, Bouncy, Crystal, Dashing, Eloquent, Foxy, Galactic, Humble, Iron, Jazzy. Cuando se nombra una distribución informalmente, basta con la primera palabra (``Humble'', ``Jazzy'').

La Tabla~\ref{tab:distros_detalle} resume las distribuciones más relevantes con su plataforma base y estado de soporte.

\begin{table}[H]
  \centering
  \caption{Distribuciones de ROS~2 y compatibilidad con Ubuntu.}
  \label{tab:distros_detalle}
  \begin{tabular}{@{}lllp{2.2cm}@{}}
    \toprule
    Distribución & Ubuntu & Tipo & Soporte \\
    \midrule
    Foxy Fitzroy        & 20.04 & LTS & Expirado (may. 2023) \\
    Galactic Geochelone & 20.04 & STD & Expirado (dic. 2022) \\
    Humble Hawksbill    & 22.04 & LTS & Hasta may. 2027 \\
    Iron Irwini         & 22.04 & STD & Expirado (nov. 2024) \\
    Jazzy Jalisco       & 24.04 & LTS & Hasta may. 2029 \\
    \bottomrule
  \end{tabular}
\end{table}

\subsection{Qué cambia entre distribuciones}

Cambiar de distribución de ROS~2 no es equivalente a actualizar un paquete Python con \texttt{pip}. Es un cambio de plataforma completo que puede afectar tres niveles distintos:

\textbf{API de rclpy / rclcpp.} Las interfaces de programación evolucionan entre distribuciones. Algunas funciones se deprecan con advertencias en una versión y se eliminan en la siguiente. El sistema de parámetros, por ejemplo, sufrió cambios significativos entre Dashing y Foxy: el método para declarar parámetros y la semántica de valores por defecto cambió de forma incompatible hacia atrás. Migrar código entre distribuciones alejadas puede requerir ajustes no triviales, especialmente en la forma de configurar QoS, crear suscriptores con opciones avanzadas, o usar características de ciclo de vida de nodos.

\textbf{Tipos de mensajes.} Los paquetes de mensajes estándar (\texttt{std\_msgs}, \texttt{geometry\_msgs}, \texttt{sensor\_msgs}, \texttt{nav\_msgs}) son generalmente estables entre distribuciones. Los mensajes de paquetes de terceros pueden cambiar sin garantías de compatibilidad. El paquete Nav2 de navegación autónoma, por ejemplo, cambió la interfaz de sus acciones entre Humble y Jazzy~\cite{macenski2020nav2}, requiriendo adaptaciones en código que dependía de él.

\textbf{Herramientas de construcción y sistema.} \texttt{colcon} es el sistema de construcción estándar de ROS~2 y es compatible entre distribuciones. Sin embargo, la versión de CMake, Python, y las herramientas de compilación de C++ cambian con Ubuntu. Código que compila sin problemas en Ubuntu 22.04 (Humble) puede fallar en Ubuntu 24.04 (Jazzy) por cambios en el compilador, cambios en las flags por defecto de seguridad, o cambios en cómo se resuelven dependencias de sistema.

\subsection{Recomendación práctica}

Para proyectos de investigación con horizonte de uno a tres años, la distribución LTS activa es siempre la elección más segura. Al momento de escribir este documento, \textbf{Humble Hawksbill} sobre Ubuntu 22.04 es la LTS con mayor ecosistema de paquetes de terceros disponibles, la mayor cantidad de documentación y ejemplos en internet, y soporte garantizado hasta mayo de 2027.

Jazzy sobre Ubuntu 24.04 es la LTS más reciente, pero aún tiene menos paquetes de terceros disponibles y algunos cambios de API (especialmente en Nav2 y en la interfaz de acciones) que requieren adaptación. Para proyectos que arrancan en 2025 o después, Jazzy es la elección natural.

Si el proyecto necesita correr en hardware embarcado (Raspberry Pi, NVIDIA Jetson, BeagleBone), la elección de distribución también está condicionada por qué versión de Ubuntu está disponible y optimizada para ese hardware. Las versiones de Jetpack de NVIDIA, por ejemplo, están atadas a versiones específicas de Ubuntu, lo que a veces fuerza a usar distribuciones de ROS~2 no-LTS o versiones antiguas, con el consiguiente compromiso en soporte a largo plazo.

La regla práctica: \textit{no actualizar de distribución a mitad de un proyecto activo salvo que haya una razón técnica imperativa}. El costo de migración es significativo y el beneficio raramente justifica la interrupción.
